\section{Discussion and Conclusion}
\label{sec:discussion}
\label{sec:conclusion}

Atom movement in zoned architectures links placement decisions across gate layers. Keeping an atom in the entanglement zone avoids storage round trips but continues to occupy a gate site and incur idle-atom excitation. Returning it to storage changes the starting point and feasibility of subsequent transport. The cost of inter-layer movement thus depends on gate sites, residency, and storage sites together.

We presented GA-LK to jointly select gate sites, resident atoms, and return storage sites at each layer boundary. Future-layer costs are estimated from the atom positions, occupancy, and clocks after each candidate completes the current layer. Feasibility checks and conflict grouping turn placement choices into executable transport. Decayed multilayer look-ahead incorporates movements required by nonadjacent interactions into current decisions.

Under the shared physical model, GA-LK improved geometric-mean fidelity over ZAC~\cite{lin2025zac} and routing-aware placement~\cite{stade2025routingaware} on both benchmark sets. It also reduced mean rearrangement batch counts and rearrangement times. Controlled comparisons showed gains from the finite-horizon configuration and genetic search. These results support jointly considering atom reuse, return locations, and subsequent transport to improve compilation quality on zoned neutral-atom architectures.
